\subsection*{Chapter 5: Implementing a GPT Model}

\paragraph{Exercise 5.1.}
With \cref{eq:gpt:params} (untied head, no QKV bias, $\vocab=\num{50257}$, $T_\text{max}=1024$):
GPT-2 medium has \num{406212608} parameters (attention \SI{100.7}{M}, FFN \SI{201.4}{M}); large has
\num{838220800} (attention \SI{236.0}{M}, FFN \SI{472.1}{M}); XL has \num{1637792000} (attention
\SI{491.6}{M}, FFN \SI{983.4}{M}). For $L=12$ the FFN holds $12(8d^2+5d)$ parameters and the
embeddings plus head hold $(2\vocab+T_\text{max})d$. Equality gives $96d+60=2\cdot\num{50257}+1024$,
i.e.\ $d\approx1057$; any width above that (in practice the next multiple of 64, $d=1088$) puts more
parameters in the FFNs than in the vocabulary-sized matrices.

\paragraph{Exercise 5.2.}
Add three keys, e.g.\ \code{drop\_emb}, \code{drop\_attn} and \code{drop\_resid}, pass
\code{drop\_attn} to \code{MultiHeadAttention} and \code{drop\_resid} to \code{drop\_shortcut}, and
use \code{drop\_emb} in \code{GPTModel}. Removing the \emph{attention-weight} dropout saves the most:
it acts on the largest activation ($BHT^2$ elements per layer), its mask must otherwise be stored or
regenerated, and many fused attention kernels have a faster code path when dropout is zero.

\paragraph{Exercise 5.3.}
The next multiple of 128 above \num{50257} is \num{50304}: 47 extra rows, $47\times768=\num{36096}$
parameters per vocabulary matrix (\num{72192} if the head is untied). Two fp32 Adam moments add
8 bytes per parameter: \SI{289}{\kilo\byte} (tied) or \SI{578}{\kilo\byte} (untied)---negligible.
During generation either slice the logits to the true vocabulary (\code{logits[..., :50257]}) or set
the padded entries to $-\infty$ before argmax/sampling, so a padded ID can never be emitted.

\paragraph{Exercise 5.4.}
From $8Nd^2=r\,(10Nd+8d^2)$ with ridge point $r$ we get $N^\star=8rd/(8d-10r)$. For the A100,
$r=312/2.039\approx153$ FLOP/byte, giving $N^\star\approx204$ tokens for $d=768$ and $\approx160$ for
$d=4096$. The second FFN GEMM, $(N\times4d)\cdot(4d\times d)$, also performs $8Nd^2$ FLOPs and moves
$2(4Nd)+2(4d^2)+2(Nd)=10Nd+8d^2$ bytes: its intensity is identical, so (ignoring tiling effects) both
GEMMs reach the ridge at the same $N^\star$.

\paragraph{Exercise 5.5.}
Without scaling, each of the $2L$ sublayers adds an update of roughly constant variance $\sigma_u^2$
to the stream, so after $\ell$ sublayers the variance is $\sigma_0^2+\ell\sigma_u^2$ and the RMS grows
like $\sqrt{\ell}$: a concave, square-root-shaped curve. Scaling the output projections by
$1/\sqrt{2L}$ shrinks each update's variance by $2L$, so the total added variance after all $2L$
sublayers is about $\sigma_u^2$ and the curve becomes nearly flat.

\paragraph{Exercise 5.6.}
Post-LN computes \code{x = norm(x + f(x))}. Without warmup, the post-LN model typically diverges (loss
spikes to NaN or plateaus at the unigram entropy) at a learning rate several times lower than the
pre-LN model tolerates, because the gradients of the last layers are large at initialization
(Xiong et al., 2020); adding a linear warmup closes much of the gap.

\paragraph{Exercise 5.7.}
(a) The explicit tanh-GELU saves, per element of the $4d$-wide hidden tensor, the input (for
\code{pow} and the outer multiply), the \code{tanh} output, the half-scaled input and the
$(1+\tanh)$ factor: about four $4d$-wide tensors. (b) \code{nn.GELU} saves only its input, removing
roughly three of them (\SI{1387}{\kilo\byte}$\to$\SI{994}{\kilo\byte} on the tiny configuration).
(c) The explicit LayerNorm saves $x-\mu$, the denominator and the normalized tensor;
\code{nn.LayerNorm} saves the input plus one mean and one reciprocal standard deviation per row, so
roughly one to two $d$-wide tensors disappear per norm.

\paragraph{Exercise 5.8.}
Causal attention computes, per head and sequence, $\sum_{i=1}^{T} i\,\dhead\cdot2 = T(T+1)\dhead$
FLOPs for $QK^\top$ and the same for $AV$, i.e.\ $2T(T+1)d$ per layer instead of $4T^2d$---a factor
$(T+1)/2T\approx\tfrac12$. For GPT-2 small the total forward FLOPs per token drop from
\num{284.8}M to \num{265.9}M at $T=1024$ (6.6\% less) and from \num{549.1}M to \num{398.1}M at
$T=8192$ (27.5\% less). If one paper counts full attention and another counts only the causal half,
their MFU numbers differ by this amount for the same hardware efficiency, increasingly so at long
context.

\paragraph{Exercise 5.9.}
Neglecting attention terms, a forward pass over $t$ tokens costs $\approx2Pt$. Without a cache,
generating $n$ tokens costs $\sum_{t=p}^{p+n-1}2Pt\approx2P(pn+n^2/2)$; with a cache, $2P(p+n)$
(prefill plus one token per step). The ratio is $(pn+n^2/2)/(p+n)$, which for $p=100$ reaches 100 at
$n^2/2=10^4$, i.e.\ $n\approx141$. Including attention makes the uncached version slightly worse still.

\paragraph{Exercise 5.10.}
Loop over sequence chunks of 128 positions: compute the chunk's logits, its cross-entropy summed over
tokens, and accumulate; wrap each chunk in \code{torch.utils.checkpoint} so the chunk logits are
recomputed in backward rather than stored. Peak logits memory falls from
$8\cdot1024\cdot\num{50257}\cdot4\,\text{B}=\SI{1.65}{\giga\byte}$ to
$8\cdot128\cdot\num{50257}\cdot4\,\text{B}=\SI{206}{\mega\byte}$ (and the same for the softmax
gradient). The price is recomputing the head GEMM in the backward pass, about $2N\vocab d$ extra
FLOPs, i.e.\ one third more work for the output layer.
